% Lösungen Einheit 4 – Nr. 50–57
\einheitenkopf[le4][4 Verdopplungs- und Halbwertszeit]{Einheit 4 von 5 · Verdopplungs- und Halbwertszeit}

\erg{50}{a) der Bestand \quad b) der Bestand \quad c) die Zeit \quad d) der Bestand \quad e) die Zeit}
\erg{51}{a) 80 \quad b) 500 \quad c) 45 \quad d) 1,8\,g \quad e) 1000\,€ (das Doppelte des Startwerts, nicht von 700\,€)}
\erg{52}{a) 1\,h \quad b) etwa 2 Jahre (Ziel 32, am nächsten liegt 36) \quad c) etwa 2\,h (Ziel 40, am nächsten 39,2) \quad d) etwa 6 Tage (Ziel 6, am nächsten 6,14)}
\erg{53}{a) etwa 3\,h (bei 200\,mg) \quad b) Ich suche 200 (die Hälfte von 400) an der senkrechten Achse, gehe waagerecht bis zum Graphen und dann senkrecht nach unten; dort lese ich die Zeit ab. \quad c) Anfangswert 60, doppelter Wert 120; waagerecht von 120 zum Graphen, senkrecht zur Zeitachse: etwa 6 Monate}
\erg{54}{a) 2 \quad b) 3 \quad c) 1 \quad d) 3 \quad e) 7 \quad f) 7 \quad g) $80 \cdot 1{,}1^7 \approx 155{,}90$; \ $500 \cdot 1{,}1^7 \approx 974{,}36$ – beide haben sich ungefähr verdoppelt, egal wo sie starten. \quad h) $1{,}1^T = 2$, \ $T = \log 2 : \log 1{,}1 \approx 7{,}27$ Jahre}
\erg{55}{a) Paul verdoppelt die Zeit statt des Bestands. Gesucht ist der Bestand 600: nach etwa 4 Jahren (622,08). \quad b) Lea hat die Achsen vertauscht. Richtig: 4\,mg auf der senkrechten Achse suchen, waagerecht zum Graphen, senkrecht nach unten und die Zeit ablesen. \quad c) Gesucht ist der Wert, der 25 am nächsten liegt: 24,89 bei $t = 4$; die Halbwertszeit ist etwa 4 Stunden.}
\erg{56}{a) In gleichen Zeiten wird immer mit demselben Faktor multipliziert; um auf das Doppelte zu kommen, braucht jeder Startwert gleich viele Schritte ($1{,}1^7 \approx 2$). \quad b) Nein: von 100 bis 200 dauert es 5 Schritte, von 200 bis 400 aber 10 Schritte.}
\erg{57}{a) 14 Uhr: 120\,mg; 20 Uhr: 60\,mg; 2 Uhr: 30\,mg \quad b) um 20 Uhr (60\,mg) \quad c) Nein: nach 24 Stunden (vier Halbwertszeiten) sind noch $240 \cdot 0{,}5^4 = 15$\,mg da.}
